\chapter{95}

\section{Mailand (IT), Dienstag, 11.
September}



«Buongiorno, Signora Ricci, mein Name ist Federico Rossi,» er deutete
auf den Polizisten neben sich, «mein Kollege Melik Sönmez.»

Sie steckten ihre Ausweise ein.

«Dürfen wir reinkommen?»

Valentina hielt sich an der Tür fest. Sie hatte die beiden erwartet. Wie
sie nun vor ihr standen, in Uniform und mit ernsten Mienen, erschien ihr
die Aktion übertrieben. War sie auf der Polizeistation zu hysterisch
aufgetreten?

«Ja, natürlich, bitte,» und damit gab sie den Weg frei und schloss
hinter ihnen die Eingangstür.

Setzen wollten sich die beiden nicht. Gut so. Valentina hatte sich
bereits vorgestellt, wie sie ihnen ein Getränk anbieten müsste. In
TV-Krimis lehnten die Polizisten zwar Angebote fast durchgehend ab. Oft
mit strengem Blick. Manchmal auch mit diesen nicht ganz glaubwürdigen
Empathiegesichtern, vor allem, wenn sie Todesnachrichten überbrachten.

Die Männer standen unterdessen in ihrem kleinen Wohnzimmer und schauten
sich suchend um. Valentina wusste es. Sie würden ihre Wohnung und die
nähere Umgebung inspizieren. Und wie die Polizistin Riva am Telefon
angekündigt hatte, sich über ein mögliches Sicherheitssystem Gedanken
machen.

«Je nachdem, wie hoch Sie gefährdet sind, wird ein Schutzprogramm
verfügt, Signora Ricci,» hatte sie hinzugefügt.

Auf Valentina hatten ihre Erklärungen wenig beruhigend gewirkt. Sie
lösten viel mehr diffuse Ängste aus. Wenn die Polizei solche Massnahmen
traf, musste sie in grösserer Gefahr sein, als ihr lieb war. Und bei
diesem Gedanken ärgerte sie sich über Louis und auch über sich selbst,
dass sie sich in diese Sache hatte hineinziehen lassen.
Sönmez kam als Erster auf Valentina zu.

«Auch hier, alles sauber,» rief ihnen Rossi vom Flur her zu, kam zurück
und stellte sich neben seinen Kollegen.

«Erstmal, Signora, schätzen wir die Gefahr eines erneuten Übergriffs bis
zum 16. als klein ein, sozusagen als unwahrscheinlich,» Sönmez schaute
Valentina über den Salontisch an, die Polizisten hatten sich unterdessen
doch gesetzt, «wir treffen heute Vorabklärungen, wie Sie wissen. Je
nachdem, wie sich die Sache entwickelt, entscheiden wir uns für
entsprechende Massnahmen. Es ist möglich, dass Sie Ihre Gewohnheiten,
Ihren Alltag verändern müssen. Im Falle höchster Gefahr, werden Sie
anderswo untergebracht, aber,» er hob beruhigend die Hand, «aktuell
bewegen wir uns auf niedrigster Gefahren-Stufe, also, machen Sie sich
nicht zu viele Sorgen. Sie können uns Tag und Nacht erreichen. Wir sind
für Sie da. Wir lassen Sie nicht hängen.»

«Danke, gut zu wissen.»

Am liebsten hätte sie geheult. Sie schaukelte sich angespannt hin und
her. Und immer wieder war da die Frage, wie es hatte geschehen können,
dass sie ungefragt eine Rolle in diesem verworrenen Krimi spielte.
Die Männer standen gleichzeitig auf. Sönmez überreichte Valentina seine
Karte.

«Mit Commissario Mantovani werden Sie morgen die weiteren Schritte
besprechen.»

Sie nickte mechanisch. Es gab keine andere Option, als das angefangene
Spiel weiterzuspielen. Die Polizisten waren immerhin freundliche Kerle.

«Sie werden morgen um neun Uhr fünfzehn von uns abgeholt. Es kommt eine
Kollegin in Zivil mit einem Privatwagen, einem weissen Fiat Punto. Ihr
Name ist Claudia Ferrari.»

Sönmez hielt Valentina ein kleines Portraitbild hin. Es zeigte eine
blondierte, junge Frau.

«Das ist sie. Sie können ihr vertrauen.»
